% ============================================================================
\clearpage
\section*{v4.9.0-r07 Source-Only: The Cone Angle and the Weak Mixing Angle}
\addcontentsline{toc}{section}{v4.9.0-r07 Source-Only: The Cone Angle and the Weak Mixing Angle}
\label{sec:v490-r07-cone-weak-mixing}
% ============================================================================

\paragraph{Purpose.}
This layer records three things.  First, the pre-registered test H14 of the relation between the cone angle of r06.2 and
the weak mixing angle, with its result.  Second, an exploration that followed it, kept as a protocol in which every
variant is counted.  Third, the structural reading that came out of the exploration, and the pre-registered hypothesis
H15 it led to.  Files, all under \texttt{llm/}:
\begin{itemize}[leftmargin=2em]
\item \texttt{Vorab\_Notiz\_H14\_Kegelwinkel\_Weinberg.md};
\item \texttt{Ergebnis\_H14\_Kegelwinkel\_Weinberg.md};
\item \texttt{Erkundungsprotokoll\_Weinberg\_Kegel.md};
\item \texttt{Vorab\_Notiz\_H15\_W\_Masse\_gamma\_L.md};
\item \texttt{Ergebnis\_H15\_W\_Masse\_gamma\_L.md};
\item scripts \texttt{k43\_test\_h14.py} to \texttt{k46\_test\_h15.py}.
\end{itemize}
Data: Particle Data Group, Review of Particle Physics
2024 with the 2025 update (summary tables and the review of the electroweak model), retrieved 2026-10-21.  Earlier layers
are not rewritten.

\noindent\textbf{Source-only status.}
This v4.9.0-r07 layer is source-only.  It does not build a PDF, does not assign a new DOI,
and does not authorize a public release.  The v4.9.0 Version DOI
\href{https://doi.org/10.5281/zenodo.23534954}{10.5281/zenodo.23534954} and the Concept DOI
\href{https://doi.org/10.5281/zenodo.19210728}{10.5281/zenodo.19210728} are unchanged by it.

\subsection*{v4.9.0-r07 test H14}
\label{subsec:v490-r07-h14}

H14 states \(\vartheta_c=2\theta_W\), equivalently \(\sin^2\theta_W=(1-\bsym)/2\).  It was formulated after the closeness
\(56.7^\circ\) against \(2\theta_W\approx57.5^\circ\) had been noticed.  The note says so, fixes the on-shell definition
\(\sin^2\theta_W=1-m_W^2/m_Z^2\) as decisive, and states the expected outcome in advance.

\paragraph{Result.}
The prediction is \(0.22533\).  The on-shell value from the masses is \(0.22321\pm0.00026\), so \(z=+8.2\).  As an exact
relation H14 is \textbf{refuted}.  The prediction lies inside the band spanned by the on-shell and the
\(\overline{\rm MS}\) value (\(0.23122\)), so as a tree-level relation it is \textbf{not decidable}.  A random simple
relation would land in that band with probability of about \(0.02\), so no outcome counts as support.

\subsection*{v4.9.0-r07 exploration, counted}
\label{subsec:v490-r07-exploration}

The exploration asked whether a physical effect can close the gap of \(+0.00212\) in \(\sin^2\theta_W\).  The gap is
equivalent to removing a contribution \(\Delta\rho_X=+0.00275\) present in nature, or to an \(m_W\) smaller by
\(110\) MeV.  The variants tried, each counted, were:
\begin{itemize}[leftmargin=2em]
\item the tribonacci value in place of \(\bsym\), which moves the prediction the wrong way;
\item removal of the top-quark contribution, which overshoots to \(0.2304\);
\item the tree-level angle from \((\alpha,G_F,m_Z)\), giving \(0.212\) or \(0.234\) depending on \(\alpha\);
\item a fraction of the top contribution, with no meaning;
\item neutrinos, eleven orders of magnitude too small;
\item every combination of omitted contributions, of which only two are large (top \(0.00934\), remainder \(0.00108\))
and none hits;
\item mixing with a dark state.
\end{itemize}
The last variant gives a general bound.  Any real contribution of the required size would shift \(m_W\) by \(155\) MeV
against the Standard-Model prediction.  The measured value differs from that prediction by \(+10\pm14\) MeV.  Real effects
of particles, known or dark, can therefore account for at most about a quarter of the gap.

\subsection*{v4.9.0-r07 the structural reading}
\label{subsec:v490-r07-structure}

\paragraph{Theoremlet v4.9.0-r07.1.}
With \(\Pi(x)=2\arctan e^{-x}\), Lobachevsky's angle of parallelism, the following are equivalent:
\[
\cos2\theta=\tanh\eta,\qquad \tan\theta=e^{-\eta},\qquad \sin^2\theta=\frac1{1+e^{2\eta}},\qquad 2\theta=\Pi(\eta).
\]
The cone angle of r06.2 is \(\vartheta_c=\Pi(\eta)\).  For the port weighting \(e^{t\sigma_x}=e^tC_{\HC}+e^{-t}A\) of
r05.1(d) at \(t=\eta\), the ratio of the weights is \(e^{2\eta}=\cot^2\theta\).

\paragraph{Proof sketch.}
\(\cos2\theta=(1-\tan^2\theta)/(1+\tan^2\theta)\) with \(\tan^2\theta=e^{-2\eta}\) gives \(\tanh\eta\);
\(\cos\Pi(x)=\tanh x\).  Checked numerically.  \(\square\)

\paragraph{Consequence.}
H14 is a statement about a ratio of couplings, \(g/g'=e^{\eta}\), the Doppler factor of way (C) in r01.  It is not a
statement about a ratio of masses.  Couplings run.  In the \(\overline{\rm MS}\) scheme \(\sin^2\theta_W\) has its minimum,
about \(0.231\), at \(\mu=M_W\), and rises above and below.  No scale reaches \(0.22533\).  The structural reading thus
points to the scheme in which the gap is larger, and it does not support the closeness to the on-shell value.

\subsection*{v4.9.0-r07 an observation, and hypothesis H15}
\label{subsec:v490-r07-h15}

\paragraph{Observation (found during the exploration).}
The on-shell coupling ratio is \(\cot\theta_W=m_W/\sqrt{m_Z^2-m_W^2}=1.86550\pm0.00140\).  The three constants of the
source in this range sit at the following distances: \(\gL=\cosh2\eta=1.86444\) at \(z=-0.8\), \(e^{\eta}=1.85419\) at
\(z=-8.1\), \(1/\bsym=1.82034\) at \(z=-32\).  By Theoremlet r01.3, \(\gL\) (way A) and \(e^\eta\) (way C) coincide at the
tribonacci point and differ by \(0.55\%\) otherwise.  The data sit at way (A).  The cone relation used way (C).

\paragraph{Why this is not evidence.}
The relation was found after the fact.  The three constants \(1.820\), \(1.854\), \(1.864\) lie in the range that
\(\cot\theta_W\) covers across schemes, \(1.823\) to \(1.866\), so some pairing is nearly expected.  The chance of a
\(1\sigma\) hit of one of the three at the on-shell value is about \(0.17\).  Reading~B contains no
\(SU(2)\times U(1)\), so the identification of the weight ratio with \((g/g')^2\) is formal.

\paragraph{Hypothesis H15 (pre-registered).}
\(\cot\theta_W=\gL\) in the on-shell definition, that is \(m_W=m_Z\gL/\sqrt{1+\gL^2}\).  Here \(\gBI\) is the root of the
state-counting equation \(\sum_{j}2\exp(-2\pi\gamma\sqrt{j(j+1)})=1\) that r03 cites, and \(\gstring=\ln2/(\pi\sqrt3)\).
Computed immediately after the note was committed and before any comparison:
\[
\gBI=0.237532957966,\qquad \gL=1.8646997646,\qquad m_W^{\rm pred}=80.3615\pm0.0018\ \text{GeV}.
\]
The comparison is with the first world average or combination of \(m_W\) published after 2026-10-21.  Known values do not
count.  The prediction lies within a few MeV of the Standard-Model prediction (\(80.359\pm0.006\) GeV).  A comparison can
therefore refute H15, but it can support H15 over the Standard Model only once the measurement excludes the
Standard-Model value at \(3\sigma\), which needs a precision below about \(1\) MeV.  The result is pending.

\subsection*{v4.9.0-r07 register entries}
\label{subsec:v490-r07-register}

\begin{center}
\small
\begin{tabular}{@{}>{\raggedright\arraybackslash}p{0.34\textwidth}>{\raggedright\arraybackslash}p{0.58\textwidth}@{}}
\toprule
statement & status after v4.9.0-r07\\
\midrule
H14, \(\vartheta_c=2\theta_W\) & refuted as exact relation (\(z=+8.2\)); not decidable at tree level; no support\\
cone angle & \(\vartheta_c=\Pi(\eta)\), the angle of parallelism (r07.1); not measured\\
gap of H14 by real particle effects & bounded by \(m_W\) to at most about a quarter (exploration, 8 counted variants)\\
H15, \(\cot\theta_W=\gL\) & pre-registered; prediction \(m_W=80.3615\pm0.0018\) GeV; comparison pending\\
r01 ways (A)/(C) & the on-shell coupling ratio lies at way (A) (observation, not a test)\\
\bottomrule
\end{tabular}
\end{center}

\subsection*{v4.9.0-r07 ledger and guard update}
\label{subsec:v490-r07-ledger-guard}

\[
R_{\rm ew}^{v490\text{-}r07}=(0,0,1)
\]
for (H14 decided, structural reading of the cone established, H15 pending).

The layer spends no guard:
\[
R_{\rm glob}^{(v490\text{-}r07,+)}=(0,1,1,1,1),
\qquad
\Gamma_{\rm guard}^{v490\text{-}r07}=(1,1,1,1,1),
\qquad
\beta_4=1/256\ \text{unspent}.
\]
It does not imply exact finite \(C5\) zero closure, exact all-readout terminality,
\(C6/T2\), a global photon-clock/RH theorem, or global all-readout closure.

